% !TEX root = ../main.tex

\section{Požiadavky na systém}

Na základe zadania boli požiadavky na riešený systém rozdelené do štyroch oblastí. Prvou oblasťou 
sú funkčné požiadavky, ktoré určujú, čo má systém umožňovať používateľovi a aké správanie sa od 
neho očakáva. Ďalšie oblasti sa týkajú hardvérovej realizácie, softvérového riadenia a dátovej 
časti vrátane IoT integrácie.

Takéto rozdelenie umožňuje oddeliť požadované funkcie systému od technických prostriedkov, 
ktorými sú tieto funkcie realizované. Funkčné požiadavky opisujú očakávané správanie systému, 
zatiaľ čo technické požiadavky určujú, aké hardvérové, softvérové a dátové časti sú potrebné 
na jeho realizáciu.

\subsection{Funkčné požiadavky}

Systém má umožňovať meranie, sledovanie a riadenie teploty kvapaliny v uzavretom hydraulickom okruhu. 
Používateľ má mať možnosť sledovať aktuálny stav systému, meniť vybrané riadiace parametre a spúšťať 
alebo zastavovať meranie.

Z funkčného hľadiska má systém zabezpečovať najmä:

\begin{itemize}
    \item meranie aktuálnej teploty chladiacej kvapaliny,
    \item reguláciu teploty pomocou Peltierovho článku a vodnej pumpy,
    \item zobrazenie aktuálnych hodnôt v reálnom čase,
    \item nastavenie požadovanej teploty,
    \item zapnutie a vypnutie systému z webového rozhrania,
    \item prepínanie medzi regulačnými módmi,
    \item manuálne nastavenie PWM hodnoty pumpy a Peltierovho článku,
    \item simuláciu tepelnej záťaže pomocou vyhrievacieho telesa,
    \item záznam priebehu merania,
    \item zobrazenie historických meraní,
    \item vzdialené monitorovanie a základné ovládanie cez ThingsBoard.
\end{itemize}

Hlavnou funkciou systému je sledovať teplotu kvapaliny a na základe nastavenej požadovanej hodnoty 
ovplyvňovať akčné členy systému. Okrem samotného merania má systém poskytovať aj používateľské rozhranie, 
cez ktoré je možné sledovať aktuálny priebeh a meniť vybrané parametre.

\subsection{Technické požiadavky na hardvér}

Hardvérová časť má vytvoriť reálnu tepelnú sústavu, na ktorej je možné pozorovať zmenu teploty 
kvapaliny a testovať rôzne spôsoby jej riadenia. Základom je uzavretý hydraulický okruh doplnený 
o výkonové prvky, senzorickú časť a riadiacu elektroniku.

Z hardvérového hľadiska je potrebné zabezpečiť:

\begin{itemize}
    \item uzavretý vodný okruh pozostávajúci z nádržky, pumpy, hadičiek a výmenníka tepla,
    \item Peltierov článok ako hlavný chladiaci prvok,
    \item účinné chladenie teplej strany Peltierovho článku pomocou chladiča a ventilátora,
    \item digitálne meranie teploty kvapaliny pomocou snímača umiestneného v nádržke,
    \item výkonové rozhranie pre riadenie Peltierovho článku, pumpy a vyhrievacieho telesa,
    \item stabilné napájanie výkonovej aj riadiacej časti systému,
    \item oddelenie výkonovej časti od riadiacej elektroniky,
    \item lokálne zobrazovanie vybraných údajov na LCD displeji,
    \item bezpečné mechanické upevnenie a prehľadné prepojenie jednotlivých častí.
\end{itemize}

Dôležitou požiadavkou je najmä bezpečná prevádzka Peltierovho článku. Teplá strana článku musí byť 
počas jeho činnosti aktívne chladená, pretože nedostatočný odvod tepla môže viesť k zníženiu účinnosti 
alebo k poškodeniu článku. Rovnako je potrebné zohľadniť vyšší prúdový odber výkonových prvkov a vhodne 
dimenzovať vodiče, svorky a napájanie.

\subsection{Požiadavky na softvér a riadenie}

Softvérová časť má zabezpečiť zber údajov z hardvéru, spracovanie nameraných hodnôt, komunikáciu 
medzi jednotlivými časťami systému a riadenie akčných členov. Riešenie je rozdelené medzi 
mikrokontrolér ESP32, backendový server a webové rozhranie.

Softvérové riešenie má zabezpečiť:

\begin{itemize}
    \item čítanie teploty zo snímača,
    \item generovanie PWM signálov pre Peltierov článok a vodnú pumpu,
    \item odosielanie nameraných údajov z mikrokontroléra na server,
    \item spracovanie aktuálneho stavu systému na backendovej strane,
    \item vytvorenie REST API pre príjem dát a ovládanie systému,
    \item realtime prenos údajov do webového rozhrania pomocou WebSocket komunikácie,
    \item webový dashboard s grafmi a ovládacími prvkami,
    \item možnosť nastavenia setpointu, regulačného módu a PWM hodnôt,
    \item prípravu regulačných algoritmov pre jednotlivé módy,
    \item podporu archivácie a zobrazenia historických meraní.
\end{itemize}

Riadenie systému má byť navrhnuté tak, aby bolo možné pracovať s tromi regulačnými módmi. 
Prvý mód využíva ako hlavný regulačný zásah výkon Peltierovho článku, druhý mód zmenu prietoku pumpy a 
tretí mód kombináciu oboch akčných členov.

\subsection{Požiadavky na dáta a IoT integráciu}

Dátová a IoT časť má zabezpečiť, aby bolo možné systém sledovať lokálne aj vzdialene a aby bolo možné 
spätne analyzovať priebeh meraní. Namerané hodnoty nemajú slúžiť iba na aktuálne zobrazenie, 
ale aj na archiváciu, testovanie regulačných módov a vyhodnocovanie správania systému.

Dátová časť má zabezpečiť:

\begin{itemize}
    \item priebežné spracovanie aktuálnych hodnôt,
    \item zobrazenie údajov vo webovom rozhraní v reálnom čase,
    \item ukladanie historických meraní do SQL databázy,
    \item záložné ukladanie meraní do súboru vo formáte JSON,
    \item možnosť načítania historických dát a ich zobrazenia v grafoch,
    \item odosielanie telemetrických údajov do platformy ThingsBoard,
    \item vytvorenie ThingsBoard dashboardu pre vizualizáciu veličín,
    \item spracovanie vybraných vzdialených príkazov z ThingsBoardu.
\end{itemize}

Do ThingsBoardu majú byť prenášané najmä veličiny potrebné na sledovanie stavu systému, 
napríklad aktuálna teplota, požadovaná teplota, regulačná odchýlka, PWM hodnota pumpy, 
PWM hodnota Peltierovho článku, aktuálny mód a informácia o tom, či je systém spustený. 
Vďaka tomu je možné systém monitorovať aj mimo lokálneho webového rozhrania.

